% Part: lambda-calculus
% Chapter: introduction
% Section: lambda-definability

\documentclass[../../../include/open-logic-section]{subfiles}

\begin{document}

\olfileid[pt]{lam}{int}{rep}
\olsection{Funções aritméticas \usetoken{S}{lambda definable}}

Como o cálculo lambda pode servir como modelo de computação? De início, nem
sequer está claro como dar sentido a essa afirmação. Para falar de
computabilidade nos números naturais, precisamos encontrar uma
representação adequada desses números. Eis uma que funciona
surpreendentemente bem.

\begin{defn}
Para cada número natural~$n$, defina o \emph{numeral de Church} $\num{n}$
como o termo lambda
\[
  \lambd[x][\lambd[y][(x(x(x(\dots x(y)))))]],
\]
onde há ao todo $n$ ocorrências de $x$.
\end{defn}

Os termos $\num{n}$ são ``iteradores'': ao receber $f$, $\num{n}$ devolve a
função que leva $y$ a $f^n(y)$. Observe que cada numeral está em forma
normal. Podemos agora dizer o que significa um termo lambda ``computar''
uma função nos números naturais.

\begin{defn}
Seja $f(x_0, \dots, x_{k-1})$ uma função parcial $n$-ária de $\Nat$ para
$\Nat$. Dizemos que um termo-$\lambd$~$F$ \emph{!!{lambda define}s}~$f$ sse,
para toda sequência de números naturais $n_0$, \dots,~$n_{k-1}$,
\[
F\, \num{n_0}\, \num{n_1} \dots \num{n_{k-1}} \red \num{f(n_0, n_1, \dots,
  n_{k-1})}
\]
se $f(n_0, \dots, n_{k-1})$ está definida, e $F, \num{n_0}\, \num{n_1}
\dots \num{n_{k-1}}$ não tem forma normal caso contrário.
\end{defn}

\begin{thm}
\ollabel{thm:lambda-def}
Uma função $f$ é parcialmente computável se, e somente se, é
!!{lambda defined} por um termo lambda.
\end{thm}

\begin{explain}
Esse teorema é um tanto surpreendente. Como modelo de computação, o cálculo
lambda é bastante simples; as únicas operações são abstração lambda e
aplicação!{} Contudo, com esses recursos escassos é possível implementar
qualquer procedimento computacional.
\end{explain}

\end{document}
